\clearpage
\VerdierPhysicalPageAnchor{0118}%
\begingroup
\small
\setlength{\abovedisplayskip}{4pt plus 2pt minus 2pt}%
\setlength{\belowdisplayskip}{4pt plus 2pt minus 2pt}%

\VerdierSemanticAnchor{ch2-1-2-6-exactness-complete}
est exacte. En prenant $W=Z$, on voit que l'application identique
$\operatorname{id}_Z\in\operatorname{Hom}_{\mathcal D}(Z,Z)$ provient d'une
application $s\colon Z\longrightarrow Y$, d'où l'existence de la section.
Cette section permet de construire un diagramme commutatif :

\VerdierSemanticAnchor{ch2-1-2-6-section-diagram}
\begin{equation*}
  \begin{tikzcd}[column sep=3.5em,row sep=3.0em]
    X \arrow[r,"{\binom{\operatorname{id}}{0}}"]
      \arrow[d,"{\operatorname{id}}"']
      & X\oplus Z \arrow[r,"{(0,\operatorname{id})}"]
        \arrow[d,"{(u,s)}"]
      & Z \arrow[d,"{\operatorname{id}}"] \\
    X \arrow[r,"u"'] & Y \arrow[r,"v"'] & Z
  \end{tikzcd}\,.
\end{equation*}

\VerdierSemanticAnchor{ch2-1-2-6-hom-five-lemma}
En appliquant le foncteur $\operatorname{Hom}_{\mathcal D}(W,\mathord\cdot)$,
l'exactitude de la suite (1.2.6.2) et le lemme des cinq entraînent que le
morphisme :

\VerdierSemanticAnchor{ch2-1-2-6-hom-isomorphism}
\[
  \operatorname{Hom}_{\mathcal D}\bigl(W,(u,s)\bigr)\colon
  \operatorname{Hom}_{\mathcal D}(W,X\oplus Z)
  \longrightarrow \operatorname{Hom}_{\mathcal D}(W,Y)
\]
est un isomorphisme. Par suite, le morphisme $(u,s)$ est un isomorphisme.

\VerdierSemanticAnchor{ch2-1-2-7}
\medskip
\noindent\textbf{Remarque 1.2.7.}\quad
Le corollaire (1.2.6) permet de montrer que tout foncteur exact
$\mathcal D'\xrightarrow{F}\mathcal D$ est additif. Tout d'abord, le triangle
$0\longrightarrow0\longrightarrow0\longrightarrow0$ est distingué dans
$\mathcal D'$ ; par suite, le triangle :

\VerdierSemanticAnchor{ch2-1-2-7-zero-triangle}
\[
  F(0)\xrightarrow{\operatorname{id}_{F(0)}}F(0)
  \xrightarrow{\operatorname{id}_{F(0)}}F(0)
  \longrightarrow F(0)[1]
\]
est distingué dans $\mathcal D$. On déduit alors de (1.2.2) que $F(0)=0$.
En appliquant le foncteur $F$ à des triangles distingués du type (1.2.6.1)
et en appliquant (1.2.6), on en déduit que $F$ est additif. Un raisonnement
analogue montre que tout foncteur cohomologique est additif.

\VerdierSemanticAnchor{ch2-1-2-8}
\medskip
\noindent\textbf{Définition 1.2.8.}\quad
Une \emph{décomposition} d'un morphisme $f\colon X\longrightarrow Y$ d'une
catégorie additive est un diagramme commutatif :

\VerdierSemanticAnchor{ch2-1-2-8-decomposition-diagram}
\begin{equation*}
  \begin{tikzcd}[ampersand replacement=\&,column sep=5.2em,row sep=3.4em]
    X \arrow[r,"f"] \arrow[d,"\wr"']
      \& Y \arrow[d,"\wr"] \\
    Z\oplus W
      \arrow[r,"{\left(\begin{smallmatrix}0&\operatorname{id}_W\\0&0\end{smallmatrix}\right)}"]
      \& W\oplus Z'
  \end{tikzcd}\,.
\end{equation*}

\VerdierSemanticAnchor{ch2-1-2-8-decomposable}
Un morphisme est dit \emph{décomposable} s'il admet une décomposition.

\VerdierSemanticAnchor{ch2-1-2-8-splitting}
Un \emph{scindage} d'un morphisme $f\colon X\longrightarrow Y$ est un
morphisme $g\colon Y\longrightarrow X$ tel que
\endgroup
